\begin{document}

\twocolumn[%
{\small\bfseries CSE558 Data Science\ \ \textbullet\ \ Midterm Project Report\ \ \textbullet\ \ Monsoon 2026\par}
\vspace{4pt}
{\LARGE\bfseries\raggedright\projecttitle\par}
\vspace{6pt}
\rule{\linewidth}{1pt}

{\small
\begin{tabularx}{\linewidth}{@{}K{2.4cm} Y@{}}
\toprule
Group name & C5XW20 \\
Dataset & kaggle.com/datasets/mlbysoham/adult-dataset (UCI Adult, archive.ics.uci.edu/dataset/2/adult) \\
Code & \href{https://github.com/durgesh-k-sharma/c5xw20-midterm}{github.com/durgesh-k-sharma/c5xw20-midterm} (scripts, README, pre-registration, outputs) \\
\bottomrule
\end{tabularx}}

\vspace{2pt}
{\small
\begin{tabularx}{\linewidth}{@{}L{0.5cm} L{3.4cm} L{1.6cm} L{4.6cm} Y@{}}
\toprule
\textbf{\#} & \textbf{Name} & \textbf{Roll no.} & \textbf{Institute email} & \textbf{Contribution} \\
\midrule
1 & Aryan Bundela & 2023141 & aryan23141@iiitd.ac.in & Dataset, size check (1); repo, README \\
2 & Durgesh Kumar Sharma & 2023214 & durgesh23214@iiitd.ac.in & Preparation, before and after (2, 3) \\
3 & Shourya Sharma & 2020129 & shourya20129@iiitd.ac.in & EDA, problem statements (4); Test 1 \\
4 & Arpit Dabas & 2023131 & arpit23131@iiitd.ac.in & Tests 2 to 4 (5); limitations (6) \\
\bottomrule
\end{tabularx}}
\vspace{12pt}
]

%-----------------------------------------------------------------------
\section{Dataset}

\field{Source:}{Kaggle, ``Adult Dataset'' (user mlbysoham). It is the file \code{adult.data}, the
training part of the UCI Adult (Census Income) data by Becker and Kohavi (1996), CC BY 4.0. UCI says
it was extracted from the 1994 US Census database. One row is one person in a survey sample. We
use \RawRows{} of the 48{,}842 UCI rows (\TrainShare\%).}

\field{What is raw about it:}{the file has no header row. Missing values are written as \code{?}:
\RawQCells{} cells in \RawQRows{} rows, in three columns. \RawDups{} rows are exact duplicates.
\code{education} and \code{education-num} carry the same information (16 levels each, one to one).
\code{capital-gain}, \code{hours-per-week} and \code{age} are capped at 99{,}999, 99 and 90. Money is in
1994 dollars. UCI already removed some records (for example age 16 or below), so the data is only lightly
dirty, and we say so rather than claim it is raw.}

{\small
\begin{tabularx}{\linewidth}{@{}L{3.0cm} Y Y@{}}
\toprule
 & \textbf{Raw file} & \textbf{After encoding} \\
\midrule
Instances (rows) & \RawRows & \CleanRows \\
Attributes (columns) & \RawCols & \CleanCols \\
Instances $\times$ attributes & \RawCells & \CleanCells \\
\bottomrule
\end{tabularx}}

\field{Size check:}{$\CleanRows\times\CleanCols=\CleanCells\ge3{,}000{,}000$, so the requirement is
met. The 108 columns are 6 numeric columns plus \OneHot{} one-hot columns (the 8 nominal columns have
$9+16+7+15+6+5+2+42=102$ levels). The response \code{income} is not counted. The margin is thin:
if we had dropped the rows with unknown values and the repeated \code{education} column, we would have
$\AltRows\times\AltCols=\AltCells$, which fails. So we keep both choices (see Section 2).}

\minihead{Key features}
The full data dictionary (type, unit, range, missingness for all 15 columns) is in the repo file \code{outputs/data\_dictionary.csv}.\par
{\footnotesize
\begin{tabularx}{\linewidth}{@{}L{2.35cm} L{1.5cm} Y@{}}
\toprule
\textbf{Feature} & \textbf{Type} & \textbf{Real-world meaning, range, missing} \\
\midrule
\code{income} & nominal & response: yearly income $>$50K or not; \RawHighShare\% are $>$50K \\
\code{sex} & nominal & Female or Male; none missing \\
\code{age} & discrete & years, 17 to 90 (90 is a cap) \\
\code{education} & nominal & highest level, 16 levels; \code{education-num} codes it 1 to 16 \\
\code{occupation} & nominal & job type, 14 levels and \code{?}; \RawQOcc{} missing \\
\code{workclass} & nominal & employer type, 8 levels and \code{?}; \RawQWC{} missing \\
\code{hours-per-week} & discrete & usual weekly hours, 1 to 99 \\
\code{capital-gain} & continuous & yearly gains in 1994 USD, 0 to 99{,}999 \\
\code{fnlwgt} & continuous & census weight: people the row stands for \\
\bottomrule
\end{tabularx}}

%-----------------------------------------------------------------------
\section{Data preparation}

{\footnotesize
\begin{tabularx}{\linewidth}{@{}L{1.9cm} L{2.3cm} Y@{}}
\toprule
\textbf{Issue found} & \textbf{What we did} & \textbf{Why, and not the alternative} \\
\midrule
No header; \code{?} for missing
 & added names; read \code{?} as a level called \code{Unknown}
 & \RawQBoth{} rows have \code{workclass} and \code{occupation} both \code{?}: one group of people with no job
   information, not random gaps. In the exploration half only \ExpUnkHigh\% of them earn $>$50K, against \ExpKnownHigh\% for the rest, so
   the group carries information. A mode fill would invent jobs. Dropping rows loses \AltDropPct\% and fails the size check. \\
\RawDups{} exact duplicate rows
 & dropped
 & they match on all 15 columns including the weight \code{fnlwgt}, so they are copies. \RawDupsNoW{} rows match when \code{fnlwgt} is ignored;
   we keep them, because they are different people with similar profiles. \\
\code{education} repeats \code{education-num}
 & kept both; tests use 4 education bands
 & dropping one column would break the size check. Bands (below HS, HS grad, some college, Bachelors or more) give strata big enough to test. \\
Capped values: \TopGain{} gains at 99{,}999, \TopHours{} at 99 hours, \TopAge{} at age 90
 & kept; three flag columns added
 & a cap is how the survey recorded a large value, not an error. Replacing it would invent a number. Flags let us leave them out. \\
IQR rule flags \IqrHours{} hours rows (\IqrHoursPct\%) and \IqrGain{} gains
 & kept all rows; used rank tests and robust standard errors
 & most people work exactly 40 hours or have zero gain, so the IQR is tiny and the rule flags normal values. Deleting them would remove real people. \\
\code{fnlwgt} survey weight
 & kept, not used in the main tests
 & course tests treat each row as one person. One weighted check for Test 1 shows if this matters. \\
8 nominal columns
 & one-hot encoded; \code{education-num} kept as a number
 & no order exists between nominal levels, so integer codes would invent one. \\
Split
 & 50/50, stratified on \code{income}, seed \SplitSeed, after removing duplicates
 & EDA used the exploration half (\NExplore{} rows); tests used the confirmation half (\NConfirm{} rows). The $>$50K share is \RateExplore\% and \RateConfirm\%. \\
\bottomrule
\end{tabularx}}

%-----------------------------------------------------------------------
\section{What changed after preparation}

{\footnotesize
\begin{tabularx}{\linewidth}{@{}L{2.9cm} L{1.9cm} Y@{}}
\toprule
\textbf{Measure} & \textbf{Before} & \textbf{After} \\
\midrule
Rows & \RawRows & \CleanRows \\
Attributes & \RawCols & \CleanCols \\
Rows $\times$ attributes & \RawCells & \CleanCells \\
Duplicate rows & \RawDups & 0 \\
\code{?} cells & \RawQCells{} (\RawQRows{} rows) & 0; \CleanUnkCells{} cells (\CleanUnkRows{} rows, \UnkRowsPct\%) are now the level \code{Unknown} \\
Capped values & not marked & flagged: \TopGain{} + \TopHours{} + \TopAge{} \\
\bottomrule
\end{tabularx}}

\egblock{Encoding took the attributes from \RawCols{} to \CleanCols, which clears the size
requirement. We lost \RawDups{} rows (\RowsLostPct\%) and nothing else. We gained a clean
\code{Unknown} level and bands for education. We did not remove any outlier, so every extreme value is
still in the data. One duplicate row held a \code{?}, so the \code{Unknown} count (\CleanUnkCells) is one lower than the raw \code{?} count (\RawQCells).}

%-----------------------------------------------------------------------
\section{Basic questions and problem statement}

All figures here use the exploration half only (\ExpN{} rows). The confirmation half was not opened until
the tests.

\begin{figure}[!ht]
  \figfile{figures/fig_income_band.pdf}{4.1cm}
  \caption{Men lead women in every education band, by \BandOneGap{} to \BandFourGap{} points; the gap
  grows with education.}
\end{figure}

\begin{figure}[!ht]
  \figfile{figures/fig_hours.pdf}{3.9cm}
  \caption{\ExpUnderW\% of women work under 40 hours a week against \ExpUnderM\% of men; mean hours are
  \ExpHoursW{} and \ExpHoursM.}
\end{figure}

\begin{figure}[!ht]
  \figfile{figures/fig_gain.pdf}{3.9cm}
  \caption{Capital gain is rare (\ExpGainShare\% of rows) and capped: all \ExpTopN{} rows at 99{,}999
  earn $>$50K, so the cap marks the response.}
\end{figure}

\begin{figure}[!ht]
  \figfile{figures/fig_unknown.pdf}{3.0cm}
  \caption{Unknown values sit in three columns; \ExpUnkBoth{} rows (\ExpUnkPct\%) lack both
  \code{workclass} and \code{occupation}.}
\end{figure}

\minihead{Basic questions}
\begin{enumerate}[label=\arabic*.]
  \item Do men earn more than women? Yes in this half: \ExpRateM\% of men and \ExpRateW\% of women are above 50K.
  \item Does more education mean higher income? Yes: the $>$50K share goes from \BandOneRate\% (below HS) to
        \BandFourRate\% (Bachelors or more); $r=\ExpCorrEdu$ with \code{education-num}.
  \item Does income rise with age? Weakly: $r=\ExpCorrAge$.
\end{enumerate}

\minihead{Our problem statements}
\noindent\fbox{\begin{minipage}{\dimexpr\linewidth-2\fboxsep-2\fboxrule\relax}
\textbf{P1 (primary).} Among people with the same occupation and education band, are women less likely
to earn above 50K than men?

\smallskip
\textbf{P2.} Is the sex gap in weekly hours different for married (civilian spouse present) people than
for never-married people?

\smallskip
\textbf{P3.} Do people with a capital gain work different hours from people with none?

\smallskip
\textbf{P4.} Is income independent of the employer type (\code{workclass})?
\end{minipage}}

\field{How they are novel:}{the basic questions compare two raw groups. P1 asks if the raw gap survives
matching on job and education. In the exploration half it does: the matched odds ratio is \ExpOneOR{}
(\ExpOneBoth{} usable strata of \ExpOneStrata), almost the same as the raw one, \ExpOneCrude. P2 asks about a
\emph{difference of two gaps}, which no single row records. P3 and P4 are closer to the basic questions. We keep them
because they stress other tests: P3 has \ExpFortyW\% of women and \ExpFortyM\% of men at exactly 40 hours, and P4 has levels with a handful of
rows.}

\field{Cause and effect:}{P1 sounds causal, but it is not. Matching on two variables leaves out age, hours,
industry and more. Occupation may also be a path through which a gap arises, and matching removes it. A causal claim would need
a design with the same people compared with only the sex changed, or a full model of how jobs are assigned. We only claim an association.}

\field{Disclosure:}{before the split, we saw two whole-data summaries (income share and mean hours by sex).
\code{PREREGISTRATION.md}, committed as \PreregCommit{} before any script ran, lists this and fixes every test below.}

%-----------------------------------------------------------------------
\section{Statistical inferences}

\field{Rows tested:}{the confirmation half, \NConfirmUsed{} people. Test 1 uses people with a known occupation
(\TOneRowsUnk{} unknown rows out) and only strata that contain both sexes, which leaves
\TOneRows{} people (\TOneNW{} women, \TOneNM{} men) in \TOneStrata{} of \TOneStrataAll{} strata (\TOneLostRows{} rows lost).
Test 2 uses \TTwoN{} people. Test 3 uses all. Test 4 drops \TFourUnk{} unknown rows.
We use $\alpha=0.05$ for every test, with a Holm correction across the four p-values.
``Matters'' means the whole 95\% interval lies beyond the smallest difference; ``does not matter'' means it lies inside; otherwise ``inconclusive''.}

%------------------------------------------------------------ Test 1
\subsection*{Test 1 (P1): do women have lower odds of income above 50K within the same job and education band?}

\tpart{a}{Hypotheses, written before running the test}
{\small
\begin{tabularx}{\linewidth}{@{}K{3.4cm} Y@{}}
\toprule
$H_0$ & $\mathrm{OR}\ge1$ \\
$H_1$ & $\mathrm{OR}<1$ \\
Significance level & $\alpha=0.05$ \\
Side & one-sided (left) \\
Test & Cochran-Mantel-Haenszel (CMH) \\
Smallest difference that matters & $\mathrm{OR}=0.80$ \\
\bottomrule
\end{tabularx}}

\egblock{OR: common odds ratio, women over men, across the occupation $\times$ education-band strata.
One-sided, because P1 says women are less likely. An OR of 0.80 means 20\% lower odds, the smallest gap we
would call worth noting.}

\tpart{b}{Other tests we considered}
\egblock{A yes-or-no outcome, a yes-or-no exposure and many strata call for CMH: it combines one $2\times2$ table per stratum.}
{\footnotesize
\begin{tabularx}{\linewidth}{@{}L{2.4cm} Y@{}}
\toprule
\textbf{Also considered} & \textbf{Why we did not use it} \\
\midrule
Plain $\chi^2$, one table & ignores the strata, so a mix of jobs could fake or hide a gap. Run anyway: $\chi^2=\TOneCrudeChi$, $p=\TOneCrudeP$, crude OR \TOneCrudeOR{} (\TOneCrudeLo{} to \TOneCrudeHi). Same decision, but not matched. \\
Logistic regression with stratum dummies & also valid, and gives OR \TOneLogOR{} (\TOneLogLo{} to \TOneLogHi), $z=\TOneLogZ$; it must drop strata with a constant outcome (\TOneLogRows{} rows, \TOneLogStrata{} strata) and fits a bigger model. CMH needs no model and is simpler. \\
Fisher exact test per stratum & counts are large, and \TOneStrata{} separate tests would need their own correction. \\
Breslow-Day test & not a test of P1: it asks if the OR is equal across strata. Tarone-adjusted $\chi^2=\TOneBD$ ($\TOneBDdf$ d.f.), $p=\TOneBDp$: mild differences, and \TOneBelow{} of \TOneDefined{} strata with a defined OR point below 1. \\
Weighted logistic (check) & with \code{fnlwgt} weights and robust errors: OR \TOneWOR{} (\TOneWLo{} to \TOneWHi). Weights do not change the story. \\
\bottomrule
\end{tabularx}}

\tpart{c}{Result}
{\small
\begin{tabularx}{\linewidth}{@{}K{3.4cm} Y@{}}
\toprule
Estimate & $\widehat{\mathrm{OR}}=\TOneOR$ (women \TOneRateW\%, men \TOneRateM\% above 50K) \\
95\% confidence interval & \TOneLo{} to \TOneHi{} (Robins-Breslow-Greenland) \\
Test statistic (d.f.) & $\chi^2_{\mathrm{CMH}}=\TOneChi$ (1); $z=\TOneZ$ \\
p-value (Holm-adjusted) & \TOneP{} (\TOneHolm) \\
Decision & reject $H_0$; it \TOneVerdict \\
\bottomrule
\end{tabularx}}

\egblock{Within the same occupation and education band, women have about three quarters lower odds of
earning above 50K (a pooled risk ratio of \TOneRR). The whole interval is below 0.80, so the gap matters.
Matching did not shrink the gap: the crude OR was \TOneCrudeOR, the matched OR is \TOneOR. Job type and
education level, as we measure them, do not explain it.}

%------------------------------------------------------------ Test 2
\subsection*{Test 2 (P2): is the sex gap in hours different for married and never-married people?}

\tpart{a}{Hypotheses, written before running the test}
{\small
\begin{tabularx}{\linewidth}{@{}K{3.4cm} Y@{}}
\toprule
$H_0$ & $D=0$ \\
$H_1$ & $D\ne0$ \\
Significance level & $\alpha=0.05$ \\
Side & two-sided \\
Test & t-test on the interaction in OLS, HC3 errors \\
Smallest difference that matters & 2 hours per week \\
\bottomrule
\end{tabularx}}

\egblock{$D$ is the men-minus-women gap in mean weekly hours among married-civ-spouse people, minus the same gap among
never-married people. Two-sided, because we did not know the direction in advance. Two hours is about one
working day a month.}

\tpart{b}{Other tests we considered}
\egblock{We compare a difference of two gaps across two factors, so we need an interaction test. We use OLS with robust errors. Cell spreads run from \TTwoSdMin{} to \TTwoSdMax{} hours, and a Brown-Forsythe test rejects equal spreads
($F=\TTwoBF$, $p=\TTwoBFp$), so we do not assume them, though the spreads are close.}
{\footnotesize
\begin{tabularx}{\linewidth}{@{}L{2.4cm} Y@{}}
\toprule
\textbf{Also considered} & \textbf{Why we did not use it} \\
\midrule
Classic OLS t-test & assumes equal variances. Run anyway: $t=\TTwoClassT$, $p=\TTwoClassP$, same decision. \\
Permutation test & makes no normality claim: sex shuffled within marital group 10{,}000 times, $p<10^{-4}$ (no shuffle reached $|D|$; the null SD is \TTwoPermSD{} against $D=\TTwoD$). Same decision. \\
Two Welch t-tests & give a gap of \TTwoWmGap{} h for married (\TTwoWmLo{} to \TTwoWmHi) and \TTwoWnGap{} h for never married (\TTwoWnLo{} to \TTwoWnHi), but no test of their difference by themselves. A Wald $z$ on the two gaps gives $z=\TTwoWaldZ$, $p=\TTwoWaldP$, the same as the OLS result. \\
Mann-Whitney, paired t-test & Mann-Whitney compares two groups, not a difference of two gaps; no person is in two groups. \\
\bottomrule
\end{tabularx}}

\tpart{c}{Result}
{\small
\begin{tabularx}{\linewidth}{@{}K{3.4cm} Y@{}}
\toprule
Estimate & $\hat D=\TTwoD$ h (gap \TTwoGapM{} h married, \TTwoGapN{} h never married) \\
95\% confidence interval & \TTwoLo{} to \TTwoHi \\
Test statistic (d.f.) & $t=\TTwoT$ (\TTwoDf) \\
p-value (Holm-adjusted) & \TTwoP{} (\TTwoHolm) \\
Decision & reject $H_0$; it \TTwoVerdict \\
\bottomrule
\end{tabularx}}

\egblock{The sex gap in hours is about \TTwoD{} hours larger among married people. The whole interval lies above 2 hours.
Mean hours: married men \TTwoMeanMM, married women \TTwoMeanMW; never-married men \TTwoMeanNM, women \TTwoMeanNW. With a standard error of \TTwoSE{} h, a test of this size can detect a difference of about \TTwoMDE{} h (2.8 standard errors, 80\% power), below our 2-hour line. This is a pattern, not a cause: we did not match the marital groups on age.}

%------------------------------------------------------------ Test 3
\subsection*{Test 3 (P3): do people with a capital gain work different hours?}

\tpart{a}{Hypotheses, written before running the test}
{\small
\begin{tabularx}{\linewidth}{@{}K{3.4cm} Y@{}}
\toprule
$H_0$ & the two hours distributions are equal (no shift) \\
$H_1$ & they differ (a shift) \\
Significance level & $\alpha=0.05$ \\
Side & two-sided \\
Test & Mann-Whitney $U$ \\
Smallest difference that matters & 2 hours per week \\
\bottomrule
\end{tabularx}}

\egblock{Groups: people with a capital gain ($n=\TThreeNG$) and with none ($n=\TThreeNN$). The estimate is the Hodges-Lehmann shift,
the median of all pairwise differences. Two-sided, because we had no direction.}

\tpart{b}{Other tests we considered}
\egblock{Hours are whole numbers, skewed, and \TThreeFortyG\% and \TThreeFortyN\% of the two groups sit at exactly 40. A rank test needs no normality.
The price is that it tests whether one group tends to have larger values, which is not the same as a gap in means.}
{\footnotesize
\begin{tabularx}{\linewidth}{@{}L{2.4cm} Y@{}}
\toprule
\textbf{Also considered} & \textbf{Why we did not use it} \\
\midrule
Welch t-test & valid at this size and answers a different question, a gap in means: \TThreeWD{} h (\TThreeWLo{} to \TThreeWHi), $t=\TThreeWT$ (\TThreeWDf), $p=\TThreeWP$. Its interval is wholly above 2 h, so it \TThreeWVerdict. \\
Brunner-Munzel test & allows unequal spreads and is the safer rank test: $|W|=\TThreeBMs$, $p=\TThreeBMp$. Same decision. \\
Two-sample z-test & needs a known $\sigma$. \\
Median test & both medians are \TThreeMedG{} h, so it has nothing to find. \\
\bottomrule
\end{tabularx}}

\tpart{c}{Result}
{\small
\begin{tabularx}{\linewidth}{@{}K{3.4cm} Y@{}}
\toprule
Estimate & Hodges-Lehmann shift \TThreeHL{} h (medians \TThreeMedG{} and \TThreeMedN{}; means \TThreeMeanG{} and \TThreeMeanN) \\
95\% confidence interval & \TThreeLo{} to \TThreeHi{} (bootstrap, 2{,}000 resamples) \\
Test statistic & $U=\TThreeU$, $z=\TThreeZ$; ($n_1=\TThreeNG$, $n_2=\TThreeNN$) \\
p-value (Holm-adjusted) & \TThreeP{} (\TThreeHolm) \\
Decision & reject $H_0$; on the Hodges-Lehmann estimate it \TThreeVerdict \\
\bottomrule
\end{tabularx}}

\egblock{The two tests disagree about size, and the reason matters. The distributions differ (a random person with a gain works more hours than a
random person without one with probability \TThreePS, counting ties as half). But the typical person in both groups works 40 hours,
so the Hodges-Lehmann shift is \TThreeHL{} and its interval has zero width. The difference is in the tails: \TThreeOverG\% of the gain group work over 40 hours against \TThreeOverN\%,
and \TThreeUnderG\% work under 40 against \TThreeUnderN\%. The pre-registered rule gives ``does not matter''; the mean gap of \TThreeWD{} hours
does matter. We report both and do not pick the one that suits us.}

%------------------------------------------------------------ Test 4
\subsection*{Test 4 (P4): is income independent of employer type?}

\tpart{a}{Hypotheses, written before running the test}
{\small
\begin{tabularx}{\linewidth}{@{}K{3.4cm} Y@{}}
\toprule
$H_0$ & income ($>$50K or not) is independent of \code{workclass} \\
$H_1$ & they are not independent \\
Significance level & $\alpha=0.05$ \\
Side & two-sided (a $\chi^2$ test has no direction) \\
Test & Pearson $\chi^2$ test of independence \\
Smallest difference that matters & Cram\'er's $V=0.10$ \\
\bottomrule
\end{tabularx}}

\egblock{\code{Unknown} rows are out. \code{Never-worked} and \code{Without-pay} merge into \code{No-pay} (\TFourNoPay{} people), because each is too small alone. $V=0.10$ is the usual
cut for a small association.}

\tpart{b}{Other tests we considered}
\egblock{A table of two categorical variables ($7\times2$) calls for a $\chi^2$ test of independence. Its rule of thumb holds: \TFourBelow{} of \TFourCells{} cells has an expected count under 5 (smallest \TFourMinE), and none is under 1.}
{\footnotesize
\begin{tabularx}{\linewidth}{@{}L{2.4cm} Y@{}}
\toprule
\textbf{Also considered} & \textbf{Why we did not use it} \\
\midrule
G-test (likelihood ratio) & gives $G=\TFourG$ ($\TFourDf$ d.f.), $p=\TFourGp$; same decision, and no better at this size. \\
Monte Carlo $\chi^2$ & 10{,}000 shuffles of income, $p<10^{-4}$. Same decision; it does not lean on the large-sample rule for the \code{No-pay} row. \\
Fisher-Freeman-Halton exact & not feasible for \TFourN{} rows; the Monte Carlo test stands in for it. \\
Z-test per level & compares only two groups; it does not test all seven together. \\
\bottomrule
\end{tabularx}}

\tpart{c}{Result}
{\small
\begin{tabularx}{\linewidth}{@{}K{3.4cm} Y@{}}
\toprule
Estimate & Cram\'er's $V=\TFourV$ (\code{Self-emp-inc} \TFourSEI\%, \code{Private} \TFourPriv\% above 50K) \\
95\% confidence interval & \TFourLo{} to \TFourHi{} (bootstrap, 2{,}000 resamples) \\
Test statistic (d.f.) & $\chi^2=\TFourChi$ (\TFourDf) \\
p-value (Holm-adjusted) & \TFourP{} (\TFourHolm) \\
Decision & reject $H_0$; it \TFourVerdict \\
\bottomrule
\end{tabularx}}

\egblock{Income depends on employer type, and the whole interval is above 0.10. The self-employed with an incorporated business have the highest $>$50K share (\TFourSEI\%); private employees have \TFourPriv\%.}

\minihead{Summary of the four tests}
{\footnotesize
\begin{tabularx}{\linewidth}{@{}L{0.5cm} L{3.4cm} L{0.8cm} L{1.5cm} Y@{}}
\toprule
\textbf{Test} & \textbf{Estimate (95\% CI)} & \textbf{Line} & \textbf{Holm p} & \textbf{Verdict} \\
\midrule
1 & OR \TOneOR{} (\TOneLo{} to \TOneHi) & 0.80 & \TOneHolm & \TOneVerdict \\
2 & $\hat D$ \TTwoD{} h (\TTwoLo{} to \TTwoHi) & 2 h & \TTwoHolm & \TTwoVerdict \\
3 & HL \TThreeHL{} h (\TThreeLo{} to \TThreeHi); mean gap \TThreeWD{} h (\TThreeWLo{} to \TThreeWHi) & 2 h & \TThreeHolm & HL: \TThreeVerdict; mean: \TThreeWVerdict \\
4 & $V$ \TFourV{} (\TFourLo{} to \TFourHi) & 0.10 & \TFourHolm & \TFourVerdict \\
\bottomrule
\end{tabularx}}

All four reject $H_0$ after the Holm correction. Only Test 3 changes its verdict with the estimate we choose.

\begin{figure}[!ht]
  \figfile{figures/fig_forest.pdf}{4.3cm}
  \caption{Each estimate in units of its smallest difference that matters. Grey is the zone where the
  difference does not matter; Test 3 lands in it only with the Hodges-Lehmann estimate.}
\end{figure}

%-----------------------------------------------------------------------
\section{Limitations}
\begin{itemize}
  \item \textbf{Old data.} The records are from 1994 and the 50K line is in 1994 dollars. The results describe that year.
  \item \textbf{Association only.} Matching leaves out age, hours and industry, and occupation may be a path of the gap, so P1 is not a proof of unequal pay.
  \item \textbf{Not fully blind.} We saw two whole-data summaries before the split (Section 4). The overall direction of P1, and of the hours gap behind P2, was already known.
  \item \textbf{Large $n$.} With about 16{,}000 rows, tiny gaps give tiny p-values, so we judge by the intervals against the thresholds. We chose those thresholds, and others could choose differently.
  \item \textbf{Test 3 depends on the estimate} (above), because of ties at 40 hours.
  \item \textbf{Independence.} The survey samples households, so spouses may both appear. The file has no household id to check this.
  \item \textbf{Weights.} Main tests ignore \code{fnlwgt}; one weighted check agrees (OR \TOneWOR). Only the training half of UCI is used.
  \item \textbf{Small cells.} \code{No-pay} has \TFourNoPay{} people; the comparison tests are not corrected for multiple testing, only the four main tests are.
\end{itemize}

\small\textbf{References.} Becker, B. and Kohavi, R. (1996). Adult. UCI Machine Learning Repository, archive.ics.uci.edu/dataset/2/adult.
Kohavi, R. (1996). Scaling up the accuracy of naive-Bayes classifiers: a decision-tree hybrid. KDD-96.

\end{document}
